\begin{proof}
\textbf{1. Norm via the inner product.}  
By definition  
\[
\|x\|=\sqrt{\langle x,x\rangle},\qquad 
\|y\|=\sqrt{\langle y,y\rangle}.
\]

\textbf{2. Degenerate cases.}  
If \(x=0\) then \(\|x+y\|=\|y\|=\|x\|+\|y\|\); similarly for \(y=0\).  
Hence we may assume \(x\neq0\) and \(y\neq0\).

\textbf{3. Squaring the inequality.}  
Since both sides are non‑negative,
\begin{align}
\|x+y\|\le\|x\|+\|y\|
\;\Longleftrightarrow\;
\|x+y\|^{2}\le(\|x\|+\|y\|)^{2}. \tag{1}
\end{align}
Using \((\|x\|+\|y\|)^{2}= \|x\|^{2}+2\|x\|\|y\|+\|y\|^{2}\) and \(\|z\|^{2}=\langle z,z\rangle\), inequality (1) becomes
\begin{align}
\langle x+y,x+y\rangle\le 
\langle x,x\rangle+2\|x\|\|y\|+\langle y,y\rangle .
\tag{2}
\end{align}

\textbf{4. Expand the left‑hand side.}  
Bilinearity and \(\langle u,v\rangle=\overline{\langle v,u\rangle}\) give
\begin{align}
\langle x+y,x+y\rangle
&=\langle x,x\rangle+\langle x,y\rangle+\langle y,x\rangle+\langle y,y\rangle \\
&=\langle x,x\rangle+2\operatorname{Re}\langle x,y\rangle+\langle y,y\rangle .
\tag{3}
\end{align}

\textbf{5. Apply Cauchy–Schwarz.}  
The Cauchy–Schwarz inequality yields
\[
|\langle x,y\rangle|\le\|x\|\,\|y\|,
\]
hence
\begin{align}
2\operatorname{Re}\langle x,y\rangle
\le 2|\langle x,y\rangle|
\le 2\|x\|\,\|y\|.
\tag{4}
\end{align}

\textbf{6. Combine (3) and (4).}  
Substituting the bound (4) into (3) we obtain exactly (2).

\textbf{7. Take square roots.}  
Both sides of (2) are non‑negative, so
\[
\|x+y\|
=\sqrt{\langle x+y,x+y\rangle}
\le\sqrt{\|x\|^{2}+2\|x\|\|y\|+\|y\|^{2}}
=\|x\|+\|y\|,
\]
which is the desired triangle inequality.

\textbf{8. Equality condition.}  
Equality can occur only where equality holds in (4), i.e. in the
Cauchy–Schwarz step.  Equality in Cauchy–Schwarz holds iff
\(y=\lambda x\) for some \(\lambda\in\mathbb{C}\) with
\(|\lambda|=\lambda\); the requirement that
\(2\operatorname{Re}\langle x,y\rangle=2\|x\|\,\|y\|\) forces \(\lambda\) to be a
non‑negative real number.  Together with the degenerate cases \(x=0\) or
\(y=0\) we obtain
\[
\|x+y\|=\|x\|+\|y\|
\iff
\bigl(x=0\bigr)\lor\bigl(y=0\bigr)\lor\bigl(y=\lambda x\ \text{with }\lambda\ge0\bigr).
\]

Thus the triangle inequality holds in every inner product space, and equality
occurs precisely when the two vectors are non‑negative scalar multiples of each
other (including the trivial zero‑vector cases). \(\square\)
\end{proof}